% ---- Macros and environments used in this solution (guarded to avoid clashes) ----
\providecommand{\bZ}{\mathbb{Z}}
\providecommand{\bQ}{\mathbb{Q}}
\providecommand{\bR}{\mathbb{R}}
\providecommand{\bN}{\mathbb{N}}
\providecommand{\norm}[1]{\lvert #1\rvert^{2}}
\providecommand{\Cut}{\operatorname{Cut}}
\providecommand{\Flow}{\operatorname{Flow}}
\providecommand{\im}{\operatorname{im}}
\providecommand{\cok}{\operatorname{cok}}
\providecommand{\cF}{\mathcal{F}}
\providecommand{\cC}{\mathcal{C}}
\makeatletter
\@ifundefined{c@thm}{\newtheorem{thm}{Theorem}}{}
\@ifundefined{c@lem}{\newtheorem{lem}{Lemma}}{}
\@ifundefined{c@cor}{\newtheorem{cor}{Corollary}}{}
\@ifundefined{c@rem}{\newtheorem{rem}{Remark}}{}
\makeatother
% ---------------------------------------------------------------------------------

\section*{Solution to Problem 6}

Throughout, $T$ is a finite weighted tree with vertex set $V$, edge set $E$, weight
function $w:V\to\bZ$, and degree function $d:V\to\bZ_{\ge 0}$.  We write the symmetric
bilinear form on $L(T)$ (and its various extensions of scalars) as $\cdot$, and we put
$\norm{x}=x\cdot x$.  Recall that a vertex $v$ is \emph{bad} if $w(v)<d(v)$, and that an
element of a positive definite lattice is \emph{irreducible} if it cannot be written as a
sum of two nonzero elements whose pairing is $\ge 0$.  We must prove:

\begin{thm}
\label{thm:one_bad_vtx}
If a weighted tree $T$ contains a single bad vertex and $L(T)$ is positive definite, then
$L(T)$ contains an irreducible vertex element.
\end{thm}

We prove the slightly sharper statement that \emph{either the unique bad vertex $v$ is
irreducible, or $T$ contains a leaf of weight $1$} (a vertex $u$ of degree $1$ and weight
$w(u)=1$ satisfies $\norm{u}=1$, the minimum nonzero value of a positive definite integral
form, and is therefore irreducible).

The strategy has two ingredients.  First we develop the cut/flow machinery of graphs and a
``projection length'' formula (\S\ref{sec:cutflow}--\ref{sec:length}).  Then we apply it to
\emph{apex trees} to show that, at a vertex of a tree with no bad vertices and weight $>$
degree, the relevant dual basis vector is irreducible
(\S\ref{sec:apex}).  Finally a local analysis at the bad vertex
(\S\ref{sec:reducible}--\ref{sec:finale}) reduces the theorem to that fact.

\section{Cut and flow lattices}
\label{sec:cutflow}

Let $G$ be a finite, loopless graph with vertex set $V(G)$ and edge set $E(G)$, possibly
with parallel edges; write $e(v,w)$ for the number of edges joining distinct vertices
$v,w$.  Fix an arbitrary orientation of the edges of $G$.  This makes $G$ a $1$--dimensional
CW complex, with integral cellular chain complex
\[
0\longrightarrow C_1(G)\overset{\partial}{\longrightarrow}C_0(G)\longrightarrow 0,
\qquad \partial(e)=v-w \text{ for an edge $e$ oriented from $w$ to $v$.}
\]
Equip $C_1(G)$ and $C_0(G)$ with the symmetric bilinear forms making $E(G)$ and $V(G)$
orthonormal bases.  Both are unimodular (in fact $\bZ^{E(G)}$ and $\bZ^{V(G)}$), hence
self-dual.  Let $\partial^\ast:C_0(G)\to C_1(G)$ be the adjoint of $\partial$, so
\[
\partial^\ast v=\sum_{e\in E(G)}(\partial e\cdot v)\,e .
\]

The \emph{cut lattice} and \emph{flow lattice} of $G$ are
\[
\Cut(G)=\im\partial^\ast\subseteq C_1(G),
\qquad
\Flow(G)=\ker\partial\subseteq C_1(G).
\]
Since $\partial^\ast$ is the adjoint of $\partial$, these are orthogonal:
$\Cut(G)\cdot\Flow(G)=0$ inside $C_1(G)$.  The images $\partial^\ast v$ ($v\in V(G)$) span
$\Cut(G)$, and a direct computation gives
\begin{equation}
\label{eq:cob}
\partial^\ast u\cdot\partial^\ast v=
\begin{cases}
d(u), & u=v,\\
-e(u,v), & u\ne v.
\end{cases}
\end{equation}
Indeed $\partial^\ast u\cdot\partial^\ast v=\sum_e(\partial e\cdot u)(\partial e\cdot v)$;
an edge with both endpoints in $\{u,v\}$ and $u\ne v$ contributes $-1$ each, while an edge
incident to $u=v$ contributes $+1$ each.

\begin{lem}[Primitivity]
\label{lem:primitive}
Both $\Flow(G)$ and $\Cut(G)$ are primitive sublattices of $C_1(G)$; that is,
$C_1(G)\cap\bigl(\Flow(G)\otimes\bQ\bigr)=\Flow(G)$ and
$C_1(G)\cap\bigl(\Cut(G)\otimes\bQ\bigr)=\Cut(G)$.
Consequently the restriction maps $C_1(G)\to\Flow(G)^\ast$ and $C_1(G)\to\Cut(G)^\ast$ are
both surjective.
\end{lem}

\begin{proof}
We may assume $G$ is connected; otherwise we argue on each connected component, over which
the chain complex splits orthogonally.

\emph{Flow.}  If $x\in C_1(G)$ and $nx\in\Flow(G)$ for some integer $n\ne 0$, then
$n\,\partial x=\partial(nx)=0$ in the torsion-free group $C_0(G)$, so $\partial x=0$ and
$x\in\Flow(G)$.  Thus $\Flow(G)$ is primitive.

\emph{Cut.}  We must show $C_1(G)\cap(\Cut(G)\otimes\bQ)=\Cut(G)$, equivalently that the
quotient $C_1(G)/\Cut(G)$ is torsion-free.  We analyze $\partial^\ast:C_0(G)\to C_1(G)$.
For connected $G$,
\[
\ker\partial^\ast=\{y\in C_0(G):\partial^\ast y=0\}
=\{y:\,y\cdot\partial w=0\ \forall w\}=(\im\partial)^\perp\cap C_0(G)=\bZ\,\mathbf 1,
\]
where $\mathbf 1=\sum_{v\in V(G)}v$: indeed $\partial^\ast y\cdot e=y\cdot\partial e$ for all
$e$, so $\partial^\ast y=0$ iff $y$ is orthogonal to $\im\partial$, and for connected $G$ the
orthogonal complement of $\im\partial$ in $C_0(G)\otimes\bQ$ is the line through $\mathbf 1$;
intersecting with $C_0(G)$ gives $\bZ\mathbf 1$ since $\mathbf 1$ is a primitive vector of
$C_0(G)=\bZ^{V(G)}$.  Hence $\partial^\ast$ factors as an injection
\[
\overline{\partial^\ast}:\ \overline C:=C_0(G)/\bZ\mathbf 1\ \hookrightarrow\ C_1(G),
\qquad \im\overline{\partial^\ast}=\Cut(G),
\]
and $\overline C$ is a free $\bZ$-module of rank $|V(G)|-1$.

Now $\Cut(G)$ is primitive in $C_1(G)$ iff $\cok\overline{\partial^\ast}=C_1(G)/\Cut(G)$ is
torsion-free, iff every invariant factor (Smith normal form entry) of $\overline{\partial^\ast}$
equals $1$.  Passing to transposes (adjoints) does not change invariant factors, so it
suffices to show that the adjoint
\[
\overline\partial:\ C_1(G)\longrightarrow \overline C^\ast,
\qquad \overline\partial=(\overline{\partial^\ast})^\ast,
\]
is \emph{surjective}.  Under the identification $C_0(G)$ self-dual, $\overline C^\ast$ is the
sublattice $(\bZ\mathbf 1)^\perp\cap C_0(G)=\ker\varepsilon$, where
$\varepsilon:C_0(G)\to\bZ$ is the augmentation $\varepsilon(\sum a_v v)=\sum a_v$; and under
this identification $\overline\partial$ is simply $\partial:C_1(G)\to\ker\varepsilon$.  But
$\partial$ \emph{is} surjective onto $\ker\varepsilon$: for connected $G$, the boundaries of
edges of any spanning tree, $\{\partial f\}$, are exactly the differences $v-v_0$ of
vertices, which generate $\ker\varepsilon$ over $\bZ$.  Therefore all invariant factors of
$\overline{\partial^\ast}$ equal $1$, $C_1(G)/\Cut(G)$ is torsion-free, and $\Cut(G)$ is
primitive.

\emph{Surjectivity of restriction.}  If $L'\subseteq M$ is a primitive sublattice of a
\emph{self-dual} lattice $M$, then a $\bZ$-basis of $L'$ extends to a $\bZ$-basis of $M$; the
dual basis of $M^\ast=M$ restricts to the dual basis of $(L')^\ast$, so the restriction map
$M=M^\ast\to(L')^\ast$ is surjective.  Since $C_1(G)$ is self-dual and $\Cut(G),\Flow(G)$ are
primitive, both restriction maps surject.
\end{proof}

\begin{rem}
\label{rem:proj}
We caution that, although each of $\Cut(G)$ and $\Flow(G)$ is primitive, the orthogonal
direct sum $\Cut(G)\oplus\Flow(G)$ is in general only a \emph{finite-index} sublattice of
$C_1(G)$ (the quotient is the critical/sandpile group).  We never use
$\Cut\oplus\Flow=C_1(G)$.  All that is needed below is the surjectivity of
$C_1(G)\twoheadrightarrow\Cut(G)^\ast$ from Lemma~\ref{lem:primitive}, together with the
orthogonality $\Cut(G)\perp\Flow(G)$.  For $x\in C_1(G)$ we write $x_C$ and $x_F$ for the
\emph{rational} orthogonal projections of $x$ onto $\Cut(G)\otimes\bQ$ and
$\Flow(G)\otimes\bQ$; then $x=x_C+x_F$ with $x_C\in\Cut(G)^\ast$ and $x_F\in\Flow(G)^\ast$
(the projection of an integral vector to the dual of a primitive sublattice of a self-dual
lattice is the restriction functional, which is integral on that sublattice).
\end{rem}

\section{The projection length formula}
\label{sec:length}

Let $G$ be a connected graph with a fixed orientation.  For an ordered partition $(A,B)$ of
$V(G)$ define the \emph{cut vector}
\[
\chi_{A,B}=\sum_{e\in E(A,B)}\varepsilon_{A,B,e}\,e\in C_1(G),
\qquad
\varepsilon_{A,B,e}=\begin{cases}+1,&e\text{ oriented }A\to B,\\-1,&e\text{ oriented }B\to A,\end{cases}
\]
where $E(A,B)$ is the set of edges with one endpoint in each part.  Then
$\chi_{A,B}=-\chi_{B,A}$, and, comparing coefficients edge by edge (edges internal to $A$
appear with opposite signs in the two terms of $\partial^\ast$ and cancel, while a cut edge
survives with the correct sign),
\begin{equation}
\label{eq:chi}
\chi_{A,B}=-\sum_{v\in A}\partial^\ast v\ \in\ \Cut(G).
\end{equation}

A \emph{two-component spanning forest} of $G$ is an unordered pair $\{A,B\}$ of vertex-disjoint
subtrees of $G$ that together contain every vertex of $G$ (equivalently, a spanning forest
with exactly two tree components).  For $x=\sum_e x_e e\in C_1(G;\bR)$ define its
\emph{square-flux} on $\{A,B\}$ by
\begin{equation}
\label{eq:flux}
\cF_{\{A,B\}}(x)=(\chi_{A,B}\cdot x)^2
=\sum_{e\in E(A,B)}x_e^2+\sum_{\substack{e,e'\in E(A,B)\\ e\ne e'}}\varepsilon_{A,B,e}\varepsilon_{A,B,e'}x_ex_{e'} .
\end{equation}
(The right-hand side is well defined: $E(A,B)$ depends only on the unordered pair, and the
products $\varepsilon_{A,B,e}\varepsilon_{A,B,e'}$ are insensitive to swapping $A$ and $B$.)

\begin{lem}
\label{lem:length}
Let $G$ be a connected graph and $x\in C_1(G)$, with orthogonal decomposition
$x=x_C+x_F$, $x_C\in\Cut(G)^\ast$, $x_F\in\Flow(G)^\ast$.  Then
\[
\norm{x_C}=\frac1\tau\sum_{\{A,B\}}\cF_{\{A,B\}}(x),
\]
the sum being over all two-component spanning forests of $G$ and $\tau$ the number of
spanning trees of $G$.
\end{lem}

\begin{proof}
A \emph{connected spanning unicyclic subgraph} $R$ is a connected spanning subgraph
containing a unique cycle $C_R$; deleting any edge of $C_R$ yields a spanning tree.  Choose
signs $\delta_{R,e}\in\{\pm1\}$ ($e\in C_R$) so that
$\omega_R=\sum_{e\in C_R}\delta_{R,e}e$ satisfies $\partial\omega_R=0$; then
$\omega_R\in\Flow(G)$ and it is unique up to sign.  Define the \emph{circulation}
\begin{equation}
\label{eq:circ}
\cC_R(x)=(\omega_R\cdot x)^2
=\sum_{e\in C_R}x_e^2+\sum_{\substack{e,e'\in C_R\\ e\ne e'}}\delta_{R,e}\delta_{R,e'}x_ex_{e'} .
\end{equation}

We record three bijections; the first two are used for the diagonal terms and the
third, together with a sign identity proved below, for the cross terms.
\begin{enumerate}[(a)]
\item For a fixed edge $e$, two-component spanning forests $\{A,B\}$ with $e\in E(A,B)$
correspond to spanning trees containing $e$.  Indeed, if $\{A,B\}$ is such a forest then
$(A\cup B)\cup\{e\}$ is connected, spanning, and has $|V(G)|-1$ edges, hence is a spanning
tree containing $e$; conversely, deleting $e$ from a spanning tree $T'\ni e$ yields a
two-component spanning forest in which $e$ is a cut edge.  The two operations are mutually
inverse.
\item Connected spanning unicyclic subgraphs $R$ with $e\in C_R$ correspond to spanning
trees \emph{not} containing $e$.  If $e\in C_R$ then $R-e$ is a spanning tree not containing
$e$ (deleting a cycle edge of a connected unicyclic graph leaves a connected spanning graph
with $|V(G)|-1$ edges); conversely, for a spanning tree $T'\not\ni e$, the graph
$T'\cup\{e\}$ is connected, spanning and has $|V(G)|$ edges, hence is unicyclic, and its
unique cycle is the fundamental cycle of $e$, which contains $e$.  Again the maps are
mutually inverse.
\item For distinct edges $e,e'$, two-component spanning forests $\{A,B\}$ with
$e,e'\in E(A,B)$ correspond to connected spanning unicyclic subgraphs $R$ with
$e,e'\in C_R$, the bijection being $R\mapsto R-\{e,e'\}$ with inverse
$\{A,B\}\mapsto (A\cup B)\cup\{e,e'\}$.  We justify below that these maps are well defined
and mutually inverse, and we then prove the sign identity
\begin{equation}
\label{eq:signid}
\varepsilon_{A,B,e}\,\varepsilon_{A,B,e'}=-\,\delta_{R,e}\,\delta_{R,e'}.
\end{equation}
\end{enumerate}

\emph{Proof that {\rm (c)} is a bijection.}
Let $R$ be a connected spanning unicyclic subgraph with $e,e'\in C_R$, and put
$F=R-\{e,e'\}$.  Then $F$ is a spanning subgraph with $|E(F)|=|V(G)|-2$ edges.  Moreover
$F$ is a forest: any cycle of $F$ would also be a cycle of $R\supseteq F$, hence equal to the
unique cycle $C_R$ of $R$; but $C_R\not\subseteq F$ since $e\in C_R\setminus E(F)$, a
contradiction.  A forest on $|V(G)|$ vertices with $|V(G)|-2$ edges has exactly two
components, each a tree, so $F$ defines a two-component spanning forest $\{A,B\}$.  We check
$e,e'\in E(A,B)$.  Suppose the two endpoints of $e$ lay in the same component, say $A$.  Then
$F\cup\{e\}$ contains a cycle through $e$; this cycle lies inside $R-e'$, so $R-e'$ would
contain a cycle, forcing $C_R\subseteq R-e'$ and hence $e'\notin C_R$, a contradiction.
Thus $e\in E(A,B)$, and symmetrically $e'\in E(A,B)$.

Conversely, let $\{A,B\}$ be a two-component spanning forest with $e,e'\in E(A,B)$, and set
$R=(A\cup B)\cup\{e,e'\}$.  Adding the cut edge $e$ to the disjoint union of trees $A$ and
$B$ joins them into a single spanning tree $T'$.  Adding the second cut edge $e'$ to $T'$
produces a connected spanning graph with $|V(G)|$ edges, hence unicyclic; its unique cycle is
the fundamental cycle of $e'$ in $T'$, and so contains $e'$.  This cycle also contains $e$:
in $T'=(A\cup B)\cup\{e\}$ the edge $e$ is the only edge crossing the cut $(A,B)$, so the
$T'$-path joining the two endpoints of $e'$ (which lie on opposite sides of the cut, since
$e'\in E(A,B)$) must traverse $e$.  Hence $e,e'\in C_R$.  The two constructions
$R\mapsto R-\{e,e'\}$ and $\{A,B\}\mapsto (A\cup B)\cup\{e,e'\}$ recover the same edge sets,
so they are mutually inverse, establishing the bijection.

\emph{Proof of the sign identity \eqref{eq:signid}.}
Fix $R$ and the corresponding forest $\{A,B\}$ as above, and choose the names of the parts so
that $\chi_{A,B}=\sum_{f\in E(A,B)}\varepsilon_{A,B,f}f$ is computed with respect to the
ordered partition $(A,B)$.  Write $\mathbf 1_B=\sum_{u\in B}u\in C_0(G)$.  We first record
that, for \emph{every} edge $f$ of $G$,
\begin{equation}
\label{eq:eps_as_boundary}
\partial f\cdot \mathbf 1_B=
\begin{cases}
\varepsilon_{A,B,f}, & f\in E(A,B),\\[2pt]
0, & f\text{ has both endpoints in }A\text{ or both in }B.
\end{cases}
\end{equation}
Indeed $\partial f=(\text{head})-(\text{tail})$, so $\partial f\cdot\mathbf 1_B$ equals
$[\text{head}\in B]-[\text{tail}\in B]$.  For a cut edge $f$ oriented from $A$ to $B$ this is
$1-0=+1=\varepsilon_{A,B,f}$, and for one oriented from $B$ to $A$ it is
$0-1=-1=\varepsilon_{A,B,f}$; for a non-cut edge both bracket-terms are equal and the value is
$0$.

Now pair the identity $\partial\omega_R=0$ with $\mathbf 1_B$:
\[
0=\partial\omega_R\cdot\mathbf 1_B
=\sum_{f\in C_R}\delta_{R,f}\,(\partial f\cdot\mathbf 1_B).
\]
Every edge of $C_R$ other than $e,e'$ lies in $F=R-\{e,e'\}=A\sqcup B$, hence has both
endpoints in the same tree (both in $A$ or both in $B$); by the second case of
\eqref{eq:eps_as_boundary} such edges contribute $0$.  By the first case of
\eqref{eq:eps_as_boundary}, the edges $e,e'\in E(A,B)$ contribute $\varepsilon_{A,B,e}$ and
$\varepsilon_{A,B,e'}$.  Thus the sum collapses to
\[
0=\delta_{R,e}\,\varepsilon_{A,B,e}+\delta_{R,e'}\,\varepsilon_{A,B,e'} .
\]
Multiplying by $\delta_{R,e}\,\varepsilon_{A,B,e'}$ and using
$\delta_{R,e}^2=\varepsilon_{A,B,e'}^2=1$ yields
$0=\varepsilon_{A,B,e}\,\varepsilon_{A,B,e'}+\delta_{R,e}\,\delta_{R,e'}$, i.e.\
$\varepsilon_{A,B,e}\,\varepsilon_{A,B,e'}=-\,\delta_{R,e}\,\delta_{R,e'}$, which is
\eqref{eq:signid}.  Both sides are unchanged if the overall sign of $\omega_R$ is flipped (it
flips both $\delta$'s) or if the parts $A,B$ are swapped (it flips both $\varepsilon$'s), so
the identity is independent of these choices.

Now sum \eqref{eq:flux} over all two-component spanning forests and \eqref{eq:circ} over
all connected spanning unicyclic subgraphs, and collect the resulting polynomial in the
variables $x_e$ monomial by monomial.

\emph{Diagonal terms $x_e^2$.}  A forest $\{A,B\}$ contributes $x_e^2$ precisely when
$e\in E(A,B)$, and a unicyclic $R$ contributes $x_e^2$ precisely when $e\in C_R$.  By
bijection (a) the number of the former equals the number $\tau_e$ of spanning trees
containing $e$, and by bijection (b) the number of the latter equals the number $\tau-\tau_e$
of spanning trees not containing $e$ (here $\tau$ is the total number of spanning trees).
Hence the total coefficient of $x_e^2$ is $\tau_e+(\tau-\tau_e)=\tau$.

\emph{Cross terms $x_ex_{e'}$, $e\ne e'$.}  In \eqref{eq:flux} each forest $\{A,B\}$ with
$e,e'\in E(A,B)$ contributes the monomial $x_ex_{e'}$ with coefficient
$2\,\varepsilon_{A,B,e}\varepsilon_{A,B,e'}$, and in \eqref{eq:circ} each $R$ with
$e,e'\in C_R$ contributes $x_ex_{e'}$ with coefficient $2\,\delta_{R,e}\delta_{R,e'}$.  By
bijection (c) these two index sets are in bijection, matching $\{A,B\}$ with
$R=(A\cup B)\cup\{e,e'\}$, and by the sign identity \eqref{eq:signid} the matched
coefficients are negatives of one another:
\[
2\,\varepsilon_{A,B,e}\varepsilon_{A,B,e'}+2\,\delta_{R,e}\delta_{R,e'}=0 .
\]
Summing over the matched pairs, the total coefficient of $x_ex_{e'}$ is $0$.

Therefore all cross terms vanish and each diagonal term has coefficient $\tau$, giving
\begin{equation}
\label{eq:full}
\sum_{\{A,B\}}\cF_{\{A,B\}}(x)+\sum_{R}\cC_R(x)=\tau\sum_e x_e^2=\tau\,\norm{x}.
\end{equation}

Now apply \eqref{eq:full} to $x_C$.  Since $x_C\perp\Flow(G)$ and $\omega_R\in\Flow(G)$, we
have $\cC_R(x_C)=(\omega_R\cdot x_C)^2=0$ for every $R$, so
$\tau\norm{x_C}=\sum_{\{A,B\}}\cF_{\{A,B\}}(x_C)$.  Finally $\chi_{A,B}\in\Cut(G)$ and
$x_F\perp\Cut(G)$ give $\chi_{A,B}\cdot x=\chi_{A,B}\cdot x_C$, whence
$\cF_{\{A,B\}}(x)=\cF_{\{A,B\}}(x_C)$ for every $\{A,B\}$.  Combining the last two
identities yields the claim.
\end{proof}

\section{Apex trees}
\label{sec:apex}

For a weighted tree $T$, the vertex set $V$ is a $\bZ$-basis of $L(T)$; let $V^\ast$ be the
dual basis of $L(T)^\ast$ and $v^\ast$ the dual of $v$.

\begin{lem}
\label{lem:apex}
Let $T$ be a weighted tree with no bad vertices, and let $v\in V$ satisfy $w(v)>d(v)$.  If
$a\in L(T)^\ast$ satisfies $a\cdot v\ne 0$, then $\norm a\ge\norm{v^\ast}$.
\end{lem}

\begin{proof}
\emph{The apex graph.}  Form the graph $G$ by adjoining to $T$ a single \emph{apex} vertex
$v_\infty$ joined to each $u\in V$ by exactly $w(u)-d(u)\ge0$ parallel edges (no edges where
$w(u)=d(u)$).  In $G$ every original vertex $u$ has degree
$d_G(u)=d(u)+(w(u)-d(u))=w(u)$, and by \eqref{eq:cob} the map $u\mapsto\partial^\ast u$ is an
isometry from $L(T)$ onto $\Cut(G)$: diagonal entries become $d_G(u)=w(u)$ and off-diagonal
entries $-e_G(u,u')=-1$ exactly when $(u,u')\in E(T)$ (there are no $T$-edges to $v_\infty$).
Thus
\[
L(T)\xrightarrow{\ \sim\ }\Cut(G),\qquad u\mapsto\partial^\ast u,
\]
and dually $L(T)^\ast\cong\Cut(G)^\ast$.  Because $w(v)>d(v)$ there is at least one edge $e$
joining $v$ to $v_\infty$.  Note $G$ is connected (it is a tree with extra apex edges, and
since $T$ is connected so is $G$).

\emph{Self-pairing of $v^\ast$.}  We claim that under
$L(T)^\ast\cong\Cut(G)^\ast$, the projection $e_C\in\Cut(G)^\ast$ of the basis edge
$e\in C_1(G)$ corresponds to $v^\ast$.  Indeed, for any $u\in V$,
\[
e_C\cdot\partial^\ast u=e\cdot\partial^\ast u
=e\cdot\sum_{f}(\partial f\cdot u)f=\partial e\cdot u,
\]
where the first equality holds because $e_C-e=-e_F\in\Flow(G)\otimes\bQ$ is orthogonal to
$\partial^\ast u\in\Cut(G)$.  Since $e$ joins $v$ to $v_\infty$, $\partial e=\pm(v-v_\infty)$,
and as $v_\infty\notin V$ we get $\partial e\cdot u=\pm\delta_{u,v}$.  After fixing the
orientation of $e$ so that the sign is $+$, this says $e_C\cdot\partial^\ast u=\delta_{u,v}$
for all $u\in V$, i.e. $e_C=v^\ast$.

Now compute $\norm{v^\ast}=\norm{e_C}$ by Lemma~\ref{lem:length} applied to $x=e$.  Here
$x_e=1$ at the single edge $e$ and $0$ elsewhere, so by \eqref{eq:flux}
\[
\cF_{\{A,B\}}(e)=
\begin{cases}1,& e\in E(A,B),\\0,&\text{otherwise.}\end{cases}
\]
By bijection (a) of Lemma~\ref{lem:length}, the number of two-component spanning forests with
$e\in E(A,B)$ equals the number $\tau_e$ of spanning trees of $G$ containing $e$.  Therefore
\begin{equation}
\label{eq:vstar}
\norm{v^\ast}=\frac{\tau_e}{\tau},
\end{equation}
where $\tau$ is the number of spanning trees of $G$.

\emph{The inequality.}  Take $a\in L(T)^\ast$ with $a\cdot v\ne0$, and let
$x_C\in\Cut(G)^\ast$ be its image under $L(T)^\ast\cong\Cut(G)^\ast$.  By the surjectivity in
Lemma~\ref{lem:primitive} there is $x\in C_1(G)$ projecting to $x_C$, i.e.
$x=x_C+x_F$.  We claim there are at least $\tau_e$ two-component spanning forests $\{A,B\}$
with $\cF_{\{A,B\}}(x)\ne0$; granting this, Lemma~\ref{lem:length} gives
\[
\norm a=\norm{x_C}=\frac1\tau\sum_{\{A,B\}}\cF_{\{A,B\}}(x)
\ \ge\ \frac{\tau_e}{\tau}=\norm{v^\ast},
\]
each nonzero square-flux being a positive integer (it equals $(\chi_{A,B}\cdot x)^2$ with
$\chi_{A,B}\cdot x\in\bZ$, because $\chi_{A,B}\in\Cut(G)\subset C_1(G)$ and $x\in C_1(G)$
both have integral coordinates).

To prove the claim we build an \emph{injection}
\[
\phi:\ \{\text{spanning trees $R$ of $G$ with $e\in R$}\}\ \longrightarrow\
\{\text{two-component spanning forests }\{A,B\}:\cF_{\{A,B\}}(x)\ne0\}.
\]

\smallskip\noindent\emph{Step 1: a useful pairing.}  Write $W=V(G)\setminus\{v\}$.  For the
connected graph $G$ one has $\sum_{u\in V(G)}\partial^\ast u=\partial^\ast(\mathbf 1)=0$,
since $\partial^\ast u\cdot e=\partial e\cdot u$ shows
$\partial^\ast(\mathbf 1)\cdot e=\partial e\cdot\mathbf 1=(\text{head}-\text{tail})\cdot
\mathbf 1=1-1=0$ for every edge $e$.  Hence, applying \eqref{eq:chi} with $A=W$,
\[
\chi_{W,\,v}=-\sum_{u\in W}\partial^\ast u
=\partial^\ast v-\sum_{u\in V(G)}\partial^\ast u=\partial^\ast v .
\]
Pairing with $x$ and using that $\partial^\ast v\in\Cut(G)$ is orthogonal to
$x_F\in\Flow(G)\otimes\bQ$, together with the isometry $\partial^\ast v\leftrightarrow v$ and
$x_C\leftrightarrow a$ under $\Cut(G)^\ast\cong L(T)^\ast$,
\begin{equation}
\label{eq:cutpair2}
\chi_{W,\,v}\cdot x=\partial^\ast v\cdot x=\partial^\ast v\cdot x_C=v\cdot a=a\cdot v\ne0 .
\end{equation}
In particular $\chi_{W,v}\cdot x\ne0$.

\smallskip\noindent\emph{Step 2: definition of $\phi$ (well-definedness).}  Fix once and for
all a linear order on the vertex set $V(G)$.  Let $R$ be a spanning tree of $G$ with $e\in R$,
and let $R_1,\dots,R_k$ be the connected components of $R-v$ (the graph $R$ with the vertex
$v$ and all $R$-edges incident to $v$ removed), indexed so that $\min R_1<\dots<\min R_k$ in
the fixed order, so that the indexing is canonical.  Because $R$ is a tree, $v$ is joined in
$R$ to each component $R_j$ by exactly one edge $g_j$ (at least one, else $R$ would be
disconnected; at most one, else $R$ would contain a cycle through $v$ and $R_j$); deleting
$g_j$ from $R$ yields exactly the two-component spanning forest $\{R_j,\,R-R_j\}$.  Thus the
forests $\{R_j,R-R_j\}$, $j=1,\dots,k$, are precisely those obtained from $R$ by deleting one
edge incident to $v$.

Since every vertex of $G-v$ lies in exactly one $R_j$, the sets $R_1,\dots,R_k$ partition
$W=V(G)\setminus\{v\}$, and \eqref{eq:chi} gives the identity of cut vectors
\[
\chi_{W,\,v}=-\sum_{u\in W}\partial^\ast u
=\sum_{j=1}^k\Big(-\sum_{u\in R_j}\partial^\ast u\Big)
=\sum_{j=1}^k\chi_{R_j,\,R-R_j}.
\]
Pairing with $x$ and invoking \eqref{eq:cutpair2},
\[
0\ne a\cdot v=\chi_{W,v}\cdot x=\sum_{j=1}^k\chi_{R_j,\,R-R_j}\cdot x,
\]
so at least one summand is nonzero.  Let $\ell$ be the \emph{smallest} index $j$ with
$\chi_{R_j,R-R_j}\cdot x\ne0$ (it exists by the previous line and is uniquely determined by
$R$ and the fixed vertex order).  Define
\[
\phi(R)=\{R_\ell,\,R-R_\ell\}.
\]
Then $\cF_{\phi(R)}(x)=(\chi_{R_\ell,R-R_\ell}\cdot x)^2\ne0$, so $\phi$ does land in the
asserted target set, and $\phi$ is a well-defined function.

\smallskip\noindent\emph{Step 3: injectivity of $\phi$ (reconstruction).}  We exhibit a map
$\psi$ defined on the image of $\phi$ with $\psi(\phi(R))=R$; this proves $\phi$ injective.
Let $\{A,B\}=\phi(R)$ be in the image, so $A$ and $B$ are vertex-disjoint subtrees of $G$
partitioning $V(G)$, and by construction
\begin{equation}
\label{eq:phistruct}
\{A,B\}=\{R_\ell,\,R-R_\ell\}\quad\text{for some component $R_\ell$ of $R-v$,}
\end{equation}
obtained from $R$ by deleting the unique edge $g_\ell$ joining $v$ to $R_\ell$.  We use only
the unordered pair $\{A,B\}$ (its vertex sets and edge sets) to recover $R$; we never use the
hidden data $\ell$ or $g_\ell$.

Note first that $v$ lies in exactly one of the two parts, namely in $R-R_\ell$ (since
$R_\ell\subseteq W=V(G)\setminus\{v\}$); call the part containing $v$ the \emph{$v$-part} and
the other the \emph{far part}.  The far part is $R_\ell$, a component of $R-v$, hence does not
contain $v$.  All of this is determined by $\{A,B\}$ alone, since "which part contains $v$"
is read off from the vertex sets.

We split into two exhaustive and mutually exclusive cases according to whether the apex edge
$e$ survives in $\{A,B\}$.  Recall $e$ joins $v$ to $v_\infty$.

\emph{Case 1: $e$ is an edge of $A$ or of $B$.}  Since $e$ joins $v$ to $v_\infty$, and the
endpoints of any edge of a part both lie in that part, $e$ being present forces both $v$ and
$v_\infty$ to lie in the same part; this must be the $v$-part.  Hence $v_\infty$ lies in the
$v$-part, so the far part is a subtree of $G$ containing neither $v$ nor $v_\infty$; in
particular the far part is a subtree of $T$ (it uses no apex edge, all of which meet
$v_\infty$).  We claim there is a \emph{unique} edge $f$ of $T$ joining $v$ to a vertex of the
far part.  At least one such edge exists: in \eqref{eq:phistruct} the deleted edge $g_\ell$
joins $v$ to $R_\ell=$ far part, and $g_\ell\ne e$ (because $e$ survives in this case, by
hypothesis), so $g_\ell$ is a non-apex, i.e.\ $T$-, edge.  At most one such edge exists: if
$f_1=(v,b_1)$ and $f_2=(v,b_2)$ were two distinct $T$-edges with $b_1,b_2$ in the (connected)
far part, then concatenating $f_1$, the path inside the far part from $b_1$ to $b_2$, and
$f_2$ would produce a cycle in the tree $T$, a contradiction.  Thus $f$ is unique and is
recoverable from $\{A,B\}$ (search $T$ for the unique edge from $v$ into the far part's vertex
set).  We set
\[
\psi(\{A,B\})=A\cup B\cup\{f\}.
\]
Since $g_\ell$ is the unique $T$-edge from $v$ to $R_\ell$ and $f$ is the unique $T$-edge from
$v$ to the far part $=R_\ell$, we have $f=g_\ell$, and therefore
$A\cup B\cup\{f\}=R_\ell\cup(R-R_\ell)\cup\{g_\ell\}=R$.

\emph{Case 2: $e$ is an edge of neither $A$ nor $B$.}  Then $e$ is the deleted edge: indeed,
the only edge of $R$ absent from $A\cup B$ is $g_\ell$, so $e\notin A\cup B$ forces $e=g_\ell$.
We set
\[
\psi(\{A,B\})=A\cup B\cup\{e\}.
\]
Then $A\cup B\cup\{e\}=R_\ell\cup(R-R_\ell)\cup\{g_\ell\}=R$.

The two cases are determined by $\{A,B\}$ alone (does $e$ appear among its edges?), are
clearly exhaustive and mutually exclusive, and in each case $\psi(\{A,B\})=R$.  Hence
$\psi\circ\phi=\mathrm{id}$ on spanning trees of $G$ containing $e$, so $\phi$ is injective.

\smallskip
Consequently the target set has at least $\tau_e$ elements, which establishes the claim and
the lemma.
\end{proof}

\begin{cor}
\label{cor:apex}
Under the hypotheses of Lemma~\ref{lem:apex}, $v^\ast\in L(T)^\ast$ is irreducible.
\end{cor}

\begin{proof}
Suppose $v^\ast=a+b$ with $a,b\in L(T)^\ast$ nonzero.  Then
$1=v^\ast\cdot v=a\cdot v+b\cdot v$, so some summand, say $a$, has $a\cdot v\ne0$.  By
Lemma~\ref{lem:apex}, $\norm a\ge\norm{v^\ast}$, while $\norm b>0$ since $b\ne0$ and the form
is positive definite.  Then
\[
\norm{v^\ast}=\norm{a+b}=\norm a+2(a\cdot b)+\norm b>\norm{v^\ast}+2(a\cdot b),
\]
forcing $a\cdot b<0$.  Hence $v^\ast$ is irreducible.
\end{proof}

\begin{cor}
\label{cor:tight}
Under the hypotheses of Lemma~\ref{lem:apex}, $0<\norm{v^\ast}\le1$, and if $\norm{v^\ast}=1$
then either $w(v)=1$ or $T$ has a leaf $u\ne v$ with $w(u)=1$.
\end{cor}

\begin{proof}
By \eqref{eq:vstar}, $\norm{v^\ast}=\tau_e/\tau$ with $0<\tau_e\le\tau$, so
$0<\norm{v^\ast}\le1$.  Moreover $\norm{v^\ast}=1\iff\tau_e=\tau\iff$ every spanning tree of
$G$ uses $e\iff e$ is a cut-edge (bridge) of $G$.  Since $G$ is $T$ together with apex edges,
$e$ (joining $v$ to $v_\infty$) is a bridge iff it is the \emph{only} apex edge of $G$, i.e.
$w(v)=d(v)+1$ and $w(u)=d(u)$ for every $u\ne v$.  If $V=\{v\}$ then $d(v)=0$ and
$w(v)=1$.  Otherwise $T$ has a leaf $u\ne v$ (a tree on $\ge2$ vertices has $\ge2$ leaves),
and for it $w(u)=d(u)=1$.
\end{proof}

\section{Reducible vertices}
\label{sec:reducible}

Let $T$ be a weighted tree with $L(T)$ positive definite, fix $v\in V$, and let
$T_1,\dots,T_k$ be the components of $T-v$.  For each $i$ let $v_i\in T_i$ be the (unique)
neighbour of $v$ in $T$.  Let $\lambda_i\in L(T_i)^\ast$ be the dual of $v_i$ with respect to
the vertex basis $V(T_i)$.  Via $L(T_i)^\ast\subset L(T_i)\otimes\bQ\subset L(T)\otimes\bQ$
we regard $\lambda_i\in L(T)\otimes\bQ$.

Set $L_1=L(T_1)\oplus\cdots\oplus L(T_k)\subset L(T)$; it is primitive, as the basis
$V\setminus\{v\}$ of $L_1$ extends to the basis $V$ of $L(T)$.  Its rational orthogonal
complement $L_1^\perp\subset L(T)\otimes\bQ$ is $1$-dimensional.  Let
$\pi:L(T)\to L_1^\perp$ be orthogonal projection and $L_0=\pi(L(T))$, a rank-$1$ rational
lattice.  Then $L(T)\subset L_0\oplus L_1^\ast=L_0\oplus L(T_1)^\ast\oplus\cdots\oplus
L(T_k)^\ast$ (orthogonal direct sum over $\bQ$).

\begin{lem}
\label{lem:irred}
$\pi(v)$ is irreducible in $L_0$.
\end{lem}

\begin{proof}
Every vertex of $T$ other than $v$ lies in some $L(T_i)\subset L_1=\ker\pi$, so projects to
$0$.  Hence $\pi(v)=\pi(\sum_{u}c_uu)$-type combinations show $\pi(v)$ generates the rank-$1$
lattice $L_0=\pi(L(T))$ (it is the image of the basis vector $v$, and all other basis vectors
die).  A generator of a rank-$1$ positive definite lattice is irreducible: if
$\pi(v)=p+q$ with $p,q$ nonzero and proportional, then $p=su,\ q=tu$ for a generator $u$ and
nonzero integers $s,t$ with $s+t=\pm1$, whence $s,t$ have opposite signs and
$p\cdot q=st\norm u<0$.
\end{proof}

\begin{lem}
\label{lem:dichotomy}
If $v$ is reducible in $L(T)$, then for some $i$ either
\textup{(i)} $\lambda_i$ is reducible in $L(T_i)^\ast$, or \textup{(ii)} $\lambda_i\in L(T_i)$.
\end{lem}

\begin{proof}
For $u\in V(T_i)$ we have $v\cdot u=-1$ if $u=v_i$ and $0$ otherwise; hence the orthogonal
projection of $v$ to $L(T_i)^\ast$ is $-\lambda_i$, and
\begin{equation}
\label{eq:vdecomp}
v=\pi(v)-\lambda_1-\cdots-\lambda_k
\end{equation}
is the orthogonal decomposition of $v$ in $L_0\oplus L(T_1)^\ast\oplus\cdots\oplus
L(T_k)^\ast$.

Suppose $v=x+y$ with $x,y\in L(T)$ nonzero and $x\cdot y\ge0$.  Decompose orthogonally
$x=\bar x+x_1+\cdots+x_k$ and $y=\bar y+y_1+\cdots+y_k$ with $\bar x,\bar y\in L_0$ and
$x_i,y_i\in L(T_i)^\ast$.  Comparing with \eqref{eq:vdecomp}, $\bar x+\bar y=\pi(v)$ and
$x_i+y_i=-\lambda_i$ for all $i$.  By Lemma~\ref{lem:irred}, $\pi(v)$ is irreducible in
$L_0$, so $\bar x\cdot\bar y\le0$, with equality only if $\bar x=0$ or $\bar y=0$.  Since
\[
0\le x\cdot y=\bar x\cdot\bar y+\sum_{i=1}^k x_i\cdot y_i,
\]
either $\bar x\cdot\bar y<0$ — in which case some $x_i\cdot y_i>0$, so the decomposition
$-\lambda_i=x_i+y_i$ exhibits $\lambda_i$ (equivalently $-\lambda_i$) as reducible in
$L(T_i)^\ast$, giving (i) — or else $\bar x\cdot\bar y=0$ and every $x_i\cdot y_i=0$.

In the latter case, equality in $\bar x\cdot\bar y\le0$ forces $\{\bar x,\bar
y\}=\{0,\pi(v)\}$, and $x_i\cdot y_i=0$ with $x_i+y_i=-\lambda_i$ irreducible (if it were
reducible we are in case (i)); thus $\{x_i,y_i\}=\{0,-\lambda_i\}$ for each $i$.  Assume
w.l.o.g. $\bar x=0$.  Then
\[
x=-\sum_{i\in A}\lambda_i\quad\text{for some }A\subseteq\{1,\dots,k\},\qquad A\ne\emptyset
\]
(as $x\ne0$).  Now $x\in L(T)$ and $x\in L_1\otimes\bQ$ (each $\lambda_i\in L(T_i)\otimes\bQ
\subseteq L_1\otimes\bQ$); since $L_1$ is primitive in $L(T)$, in fact $x\in L_1=
\bigoplus_iL(T_i)$.  As the $\lambda_i$ lie in pairwise orthogonal summands, this forces
$\lambda_i\in L(T_i)$ for every $i\in A\ne\emptyset$, which is (ii).
\end{proof}

\section{Proof of the theorem}
\label{sec:finale}

\begin{proof}[Proof of Theorem~\ref{thm:one_bad_vtx}]
Let $v$ be the unique bad vertex of $T$.  Define $T_1,\dots,T_k$, $v_i$, $\lambda_i$ as in
\S\ref{sec:reducible}.

Throughout, each subtree $T_i$ is given the weight function inherited from $T$ by restriction
(the weight $w(u)$ of a vertex $u\in V(T_i)$ is unchanged), so that "bad in $T_i$" refers to
the comparison of $w(u)$ with the degree $d_{T_i}(u)$ \emph{computed inside $T_i$}.

\smallskip\noindent\emph{Degree monotonicity under deletion of $v$.}
Removing the vertex $v$ (together with the edges incident to it) from $T$ destroys exactly the
$k$ edges $v v_i$, one into each component $T_i$.  Consequently, for every $u\in V(T_i)$,
\begin{equation}
\label{eq:degmono}
d_{T_i}(u)=\begin{cases} d_T(v_i)-1, & u=v_i,\\[2pt] d_T(u), & u\ne v_i,\end{cases}
\qquad\text{so in all cases}\quad d_{T_i}(u)\le d_T(u).
\end{equation}
(The neighbour $v_i$ loses precisely its edge to $v$; no other vertex of $T_i$ is adjacent to
$v$, so no other degree changes.)

\smallskip\noindent\emph{Badness monotonicity: $T_i$ has no bad vertex.}
Let $u\in V(T_i)$.  Since $v$ is the \emph{unique} bad vertex of $T$ and $u\ne v$ (as
$u\in T_i\subseteq T-v$), $u$ is not bad in $T$, i.e. $w(u)\ge d_T(u)$.  Combining with the
monotonicity \eqref{eq:degmono},
\[
w(u)\ \ge\ d_T(u)\ \ge\ d_{T_i}(u),
\]
so $u$ is not bad in $T_i$.  As $u\in V(T_i)$ was arbitrary, \emph{$T_i$ contains no bad
vertex.}  (Intuitively: deletion of $v$ only \emph{lowers} degrees while leaving weights
fixed, so it can only turn bad vertices good, never the reverse; and the one bad vertex $v$
of $T$ is itself removed.)

\smallskip\noindent\emph{Strict inequality at the apex vertex $v_i$.}
The vertex $v_i\ne v$ is not bad in $T$, so $w(v_i)\ge d_T(v_i)$.  Using the first line of
\eqref{eq:degmono}, $d_T(v_i)=d_{T_i}(v_i)+1$, whence
\begin{equation}
\label{eq:strictvi}
w(v_i)\ \ge\ d_T(v_i)\ =\ d_{T_i}(v_i)+1\ >\ d_{T_i}(v_i).
\end{equation}
Thus $w(v_i)>d_{T_i}(v_i)$ \emph{strictly}.

\smallskip\noindent\emph{Applying the apex results to $(T_i,v_i)$.}
The two displayed facts---$T_i$ has no bad vertex, and $w(v_i)>d_{T_i}(v_i)$---are exactly the
hypotheses of Lemma~\ref{lem:apex} (and hence of Corollaries~\ref{cor:apex} and~\ref{cor:tight},
which assume nothing more).  Note that $\lambda_i$, defined as the dual of $v_i$ with respect
to the vertex basis $V(T_i)$ of $L(T_i)$, is precisely the dual basis vector
$v_i^\ast\in L(T_i)^\ast$, so $\norm{\lambda_i}=\norm{v_i^\ast}$.  Therefore
Corollary~\ref{cor:apex} applies and shows that
\[
\lambda_i\ \text{is irreducible in}\ L(T_i)^\ast\qquad\text{for every }i=1,\dots,k.
\]
In particular case \textup{(i)} of the dichotomy Lemma~\ref{lem:dichotomy} (some $\lambda_i$
reducible in $L(T_i)^\ast$) cannot occur.

\emph{Case 1: $\lambda_i\notin L(T_i)$ for all $i$.}  Since no $\lambda_i$ is reducible (just
shown) and no $\lambda_i$ lies in $L(T_i)$, Lemma~\ref{lem:dichotomy} (contrapositive) shows
$v$ is \emph{irreducible} in $L(T)$.  Then $v$ itself is the required irreducible vertex.

\emph{Case 2: $\lambda_i\in L(T_i)$ for some $i$.}  Then $\norm{\lambda_i}\in\bZ$.  By
Corollary~\ref{cor:tight} applied to $T_i$ and $v_i$ (whose hypotheses we verified above, and
noting $\lambda_i=v_i^\ast$ in $L(T_i)^\ast$ so $\norm{\lambda_i}=\norm{v_i^\ast}$), we have
$0<\norm{\lambda_i}\le1$; integrality forces $\norm{\lambda_i}=1$, and then either
$w(v_i)=1$ or $T_i$ has a leaf $u\ne v_i$ with $w(u)=1$.

If $w(v_i)=1$: then $d_{T_i}(v_i)\le w(v_i)=1$, and since $v_i$ has a neighbour in $T_i$
(unless $T_i=\{v_i\}$) we examine its degree in $T$.  In $T$, $v_i$ is adjacent to $v$, so
$d_T(v_i)=d_{T_i}(v_i)+1$.  As $v_i$ is not bad, $w(v_i)\ge d_T(v_i)=d_{T_i}(v_i)+1$, i.e.
$1\ge d_{T_i}(v_i)+1$, forcing $d_{T_i}(v_i)=0$.  Hence $T_i=\{v_i\}$ and $v_i$ is a leaf of
$T$ (its only $T$-neighbour is $v$) with $w(v_i)=1$.

If instead $T_i$ has a leaf $u\ne v_i$ with $w(u)=1$: a leaf of $T_i$ different from $v_i$ has
the same neighbourhood in $T$ as in $T_i$ (only $v_i$ is affected by removing $v$), so $u$ is
also a leaf of $T$, with $w(u)=1$.

In either subcase, $T$ contains a leaf $u$ with $w(u)=1$, i.e. $\norm u=w(u)=1$.  Such $u$ has
minimal nonzero self-pairing in the positive definite integral lattice $L(T)$, hence is
irreducible: if $u=a+b$ with $a,b$ nonzero then $\norm a,\norm b\ge1$ and
$1=\norm u=\norm a+2(a\cdot b)+\norm b\ge2+2(a\cdot b)$, giving $a\cdot b\le-\tfrac12<0$.

In all cases $L(T)$ contains an irreducible vertex element, completing the proof.
\end{proof}

\begin{rem}[Sharpness]
The conclusion cannot be strengthened to ``the bad vertex is irreducible.''  Let $T$ be the
star with centre $v$ ($d(v)=3$, $w(v)=2$) and three leaves of weights $1,2,3$.  Then $v$ is
the unique bad vertex, $L(T)$ is positive definite of determinant $1$ (indeed isometric to
$\bZ^4$ with the standard form), yet $v=(-u_1)+(v+u_1)$, where $u_1$ is the weight-$1$ leaf,
realizes $v$ as reducible (the two summands pair to $0$).  Consistently with the theorem, the
weight-$1$ leaf $u_1$ is irreducible.
\end{rem}
